\section{Local transfer matrices and fusion}\label{sec:local}

We first describe the operator for one tile and prove its two rank reductions, its inverse relation, and its interchange identity. The resulting row identities will let us change the number and order of site parameters while keeping control of the transfer matrix.

The local algebra belongs to the integrable dilute-loop tradition. Grimm and Pearce~\cite{GrimmPearce1993} give a braid--monoid construction of dilute Yang--Baxter operators and explain their gauge equivalence with the Izergin--Korepin vertex model. We specify the gauge and scalar factors used here and verify the identities for the actual cap contractions.

\subsection{The local operator}

Set \(l=\pi/8,\ x=(2\cos l)^{-1},\ q=e^{2il},\ p=q^{-1}\). Let \(V\) have basis \(|m\rangle\), \(m=0,+1,-1\), with \(s|m\rangle=m|m\rangle\). The state \(0\) is vacant and the states \(\pm1\) are occupied. All transposes below are taken in these bases, without complex conjugation. Put \(C=\sum_{m=\pm1}p^m|m,-m\rangle\) and \(P|mk\rangle=|km\rangle\), and define the five scalar weights by
{\small
\[
 (a,b,h,d,e)(t)=
 \frac{(\sin2l\sin(3l-t),\,\sin2l\sin t,\,\sin(3l-t)\sin t,\,
 \sin(3l-t)\sin(2l-t),\,\sin(t-l)\sin t)}{\sin(2l+t)\sin(3l+t)} .
\]
}
Define \(R(t):V^{\otimes2}\to V^{\otimes2}\) by the following rules, where \(m,k=\pm1\) and the scalar weights have argument \(t\):
\[
 \begin{split}
 |00\rangle&\mapsto |00\rangle+b C,\qquad
 |m0\rangle\mapsto h|m0\rangle+a|0m\rangle,\quad |0m\rangle\mapsto h|0m\rangle+a|m0\rangle,\\
 |mk\rangle&\mapsto d|km\rangle+{\bf1}_{m=-k}p^{-m}(b|00\rangle+eC).
 \end{split}
\]
The rules conserve \(s_1+s_2\). Write \(\check R=R P\).

\subsection{Rank reductions and inversion}

At the two arguments \(2l\) and \(3l\), the range of \(R\) comes from a smaller space. For \(r=2\), let \(E=E_2:V\to V^{\otimes 2}\) take \(|0\rangle\) to \(|00\rangle+xC\) and \(|m\rangle\) to \(x(|m0\rangle+|0m\rangle)\) for \(m=\pm1\). For \(r=3\), let \(E=E_3:\mathbb C\to V^{\otimes 2}\) take 1 to \(|00\rangle+C\). The rank reduction and its compatible action on a third site are
\[
 R(rl)=EE^t P,\quad
 R_{23}(t)R_{13}(t+rl)(E\otimes I)=(E\otimes I)S,\quad S=
 \begin{cases}R(t+l)&r=2,\\I&r=3.\end{cases}
\]
Here \(S\) acts on the source of \(E\) and the third site; for \(r=3\) it is the identity on that third site. We also have the inverse relation \(R_{12}(t)R_{21}(-t)=I\). All these are meromorphic identities in their arguments.

To verify the rank reductions, substitute \(t=rl\) in the five weights and use \(C^t C=0\). We give the coefficient checks for the compatible action, since its scalar normalization will be used later. Use subscripts \(0,j,*\) for the arguments \(t,t+rl,t+l\). For \(r=2\), expand both sides on the inputs \(|00\rangle,|m0\rangle,|0m\rangle,|mk\rangle\) of \(V\otimes V\). The resulting coefficient equalities are grouped in that order in the following table:
\[
\begin{array}{ll}
 b_0+x a_jd_0=x b_*,& b_j a_0+x h_j h_0=x,\quad b_j h_0+x h_j a_0=x b_*;\\
 h_j+b_j b_0=h_*,&h_j b_0+b_j e_0=0,\quad a_j a_0+h_0=h_*,\\
 x(a_jh_0+a_0)=a_*,&b_j d_0=a_*;\\
 h_j h_0+x b_j a_0=h_*,& h_j a_0+x b_j h_0=x a_*,\quad a_j+x e_j b_0=x a_*,\\
 e_j d_0=h_*,& a_j b_0+x d_j d_0+x e_j e_0=0;\\
 d_j h_0+a_j a_0=d_*,&d_j a_0+a_j h_0=0,\quad h_j d_0=d_*,\\
 x(h_j b_0+b_j)=b_*,& b_j b_0+h_j e_0=e_*,\quad e_j a_0=b_*,\quad e_j h_0=e_* .
\end{array}
\]
The four groups occupy \(1,2,2,2\) rows. For example, the first group compares the three output types \(\sum p^m|0m,-m\rangle,\ \sum p^m|m,-m,0\rangle,\ \sum p^m|m0,-m\rangle\) on input \(|00\rangle\); the extra cap terms vanish because \(p^2+p^{-2}=0\). Each scalar equality follows on inserting the fractions and using \(2\sin A\sin B=\cos(A-B)-\cos(A+B)\) and \(8l=\pi\). For \(r=3\), the same substitution leaves the first and third groups, with the explicit \(x=h_*=1,\ a_*=b_*=0\); these are the required identities when site 3 is vacant and occupied, respectively.

The inverse relation reduces in the same way to
\[
 dd^-=1,\quad aa^-+hh^-=1,\quad ah^-+ha^-=0,\quad b+d b^-=0,\quad
 b b^-+d e^-+e d^-=0,
\]
where a superscript minus denotes the argument \(-t\). Direct substitution proves all five equalities.

\subsection{Interchange and row identities}

We next derive the interchange identity from the two rank reductions. Transpose the fusion relation and then conjugate by the operation reversing all spins. This combined operation leaves \(R\) unchanged, while spin reversal takes \(E\) to \(PE\). It therefore gives the opposite-product intertwining with \(E^t P\). Combining the two intertwinings with \(R(rl)=EE^tP\) proves
\[
 R_{12}(u-v)R_{13}(u)R_{23}(v)=R_{23}(v)R_{13}(u)R_{12}(u-v)
\]
when \(u-v=rl\), \(r=2,3\). It also holds at \(u-v=0\), since \(R(0)=P\). Multiplication by the inverse for the \(23\) pair identifies this equation with the \(132\) relation at differences \(u,u-v,-v\). The same three special cases therefore prove it when \(u=0,2l,3l\).

These six specializations determine the identity for every \(u\). Indeed, after clearing denominators, the powers of \(e^{iu}\) in each entry range from \(-4\) to \(4\) and have one parity: changing \(u\) by \(\pi\) conjugates by the sign of occupation at site 1. After removing a monomial factor, an entry is therefore a polynomial of degree at most four in \(e^{2iu}\). For generic \(v\), the six specializations are distinct modulo \(\pi\), so interpolation makes every entry zero. Meromorphic continuation gives the asserted identity for all parameters.

For a list of labelled site parameters \(\mathbf b=(b_1,\ldots,b_N)\) and \(\zeta\in\mathbb C^*\), define the spin row transfer by
\[
 T^\zeta(u;\mathbf b)=\operatorname{Tr}_a\big(\zeta^{s_a} R_{a1}(u-b_1)\cdots R_{aN}(u-b_N)\big).
\]
For a fixed twist \(\zeta\), these operators commute at all regular row arguments. To see this for \(N>0\), conjugate the trace over two auxiliary spaces by their \(R\), then use interchange at each site and conservation of the auxiliary spins. The same assertion holds for the empty list with the evident convention. Interchange also gives the adjacent exchange relation
\[
 T^\zeta(u;\ldots,A,B,\ldots)\check R_{j,j+1}(A-B)
 =\check R_{j,j+1}(A-B)T^\zeta(u;\ldots,B,A,\ldots).
\]
The analogous cyclic operation is as follows. Rotate the tensor factors one place to the left and multiply by \(\zeta^{s_1}\) on the site sent to the end. This map intertwines the row with the cyclically rotated list, by conservation in slot 1 and cyclicity of the auxiliary trace for entries at separate sites. Shifting a site parameter by \(\pi\) conjugates the row by the sign of occupation at that site.

Finally, specialize adjacent entries to \(d_0+rl,d_0\). The row intertwines through \(E_r\) with the shorter list obtained by replacing this pair with \(d_0+l\) for \(r=2\), or by deleting it for \(r=3\). This follows by inverting the fusion identity with site 3 as the auxiliary space and \(t=d_0-u\). The row relations, like the local ones, are meromorphic identities in all parameters.
